% Shared diagram for the reader, Lecture 3 and the TIM export.
% Arrows point from a foreign-key attribute to the referenced primary key.
\begin{tikzpicture}[
  x=1cm,y=1cm,
  attr/.style={anchor=west,font=\ttfamily\fontsize{10.5}{12}\selectfont,inner sep=0pt},
  title/.style={font=\sffamily\bfseries\fontsize{11}{13}\selectfont,inner sep=0pt},
  fk/.style={-{Latex[length=2.5mm,width=1.8mm]},line width=1pt,draw=blue!45!black},
  fklabel/.style={font=\sffamily\fontsize{9}{11}\selectfont,fill=white,inner sep=2pt},
  note/.style={anchor=north west,align=left,text width=4.15cm,
    font=\sffamily\fontsize{9.5}{12}\selectfont,inner sep=0pt}
]
% DELIVERY: its first three attributes form ONE primary key.
\fill[blue!7] (0,0) rectangle (4.2,-0.6);
\draw[line width=0.7pt] (0,0) rectangle (4.2,-3.15);
\draw[line width=0.7pt] (0,-0.6) -- (4.2,-0.6);
\node[title] at (2.1,-0.3) {DELIVERY};
\node[attr] at (0.22,-0.98) {\underline{supplier\_no}};
\node[attr] at (0.22,-1.58) {\underline{part\_no}};
\node[attr] at (0.22,-2.18) {\underline{project\_no}};
\node[attr] at (0.22,-2.78) {quantity};

% Referenced tables, each with its own primary key.
\fill[black!7] (7,1.6) rectangle (11.2,1.0);
\draw[line width=0.7pt] (7,1.6) rectangle (11.2,-0.95);
\draw[line width=0.7pt] (7,1.0) -- (11.2,1.0);
\node[title] at (9.1,1.3) {SUPPLIER};
\node[attr] at (7.22,0.62) {\underline{supplier\_no}};
\node[attr] at (7.22,0.02) {name};
\node[attr] at (7.22,-0.58) {city};

\fill[black!7] (7,-1.4) rectangle (11.2,-2.0);
\draw[line width=0.7pt] (7,-1.4) rectangle (11.2,-3.35);
\draw[line width=0.7pt] (7,-2.0) -- (11.2,-2.0);
\node[title] at (9.1,-1.7) {PART};
\node[attr] at (7.22,-2.38) {\underline{part\_no}};
\node[attr] at (7.22,-2.98) {name};

\fill[black!7] (7,-3.8) rectangle (11.2,-4.4);
\draw[line width=0.7pt] (7,-3.8) rectangle (11.2,-6.35);
\draw[line width=0.7pt] (7,-4.4) -- (11.2,-4.4);
\node[title] at (9.1,-4.1) {PROJECT};
\node[attr] at (7.22,-4.78) {\underline{project\_no}};
\node[attr] at (7.22,-5.38) {name};
\node[attr] at (7.22,-5.98) {budget};

% Three separate references, with no crossing arrows.
\draw[fk] (4.2,-0.98) -- (5.3,-0.98) -- (5.3,0.62)
  -- node[fklabel,above] {FK 1} (7,0.62);
\draw[fk] (4.2,-1.58) -- (5.9,-1.58) -- (5.9,-2.38)
  -- node[fklabel,above] {FK 2} (7,-2.38);
\draw[fk] (4.2,-2.18) -- (5.0,-2.18) -- (5.0,-4.78)
  -- node[fklabel,above] {FK 3} (7,-4.78);

\node[note] at (0,-3.65) {Primary keys are\\underlined.};
\node[note] at (0,-4.65) {In DELIVERY, the three\\underlined attributes form\\\textbf{one composite}\\\textbf{primary key}.};
\end{tikzpicture}
